\subsection{Engineering Validation and System Reliability}\label{sec:engineering}

A contract is only as trustworthy as the tests that execute it, and assertion
counts are not evidence. We therefore audited the suite along three
count-independent axes: mutation testing (can it detect injected semantic
faults?), assertion-strength grading (do assertions discriminate, or are they
tautological or self-referential?), and measured coverage.

\paragraph{Mutation testing: is the contract alive?} A contract whose tests cannot
fail is a document, not an executable artefact; the audit therefore injected 44
single-fault semantic mutants---none syntactic or import errors, all loading and
executing---distributed over the JavaScript plugin (29), the Python engine (11),
and the rule file (4). Each mutant was applied alone and the \emph{complete}
suites were run: 1{,}249 assertions per round, about 55{,}000 assertion
executions. A mutant is counted as killed only when an assertion fails, never
because the process crashes; no mutant was error-killed. The suite as found killed
$38/44=86.4\%$ (JavaScript $24/29=82.8\%$, Python $10/11=90.9\%$, rules
$4/4=100\%$). Of the 40 code mutants, 34 were killed by behavioural assertions
($85.0\%$); the four rule mutants were caught by cross-file consistency
assertions rather than by behaviour, a distinction we keep visible. After the four
blind spots below were pinned by minimal regression assertions, the suite kills
$42/44=95.5\%$ and $42/42=100\%$ of non-equivalent mutants. That the score moves
from 86.4\% to 95.5\% by adding tests is the concrete sense in which this contract
is alive rather than merely present.

\paragraph{Adversarial procedure.} The audit was adversarial by construction.
Mutation testing was run by one reviewer, assertion and coverage analysis by a
second, and a subset re-run independently by a third on a separate copy. All
work happened on a copy; the 30 implementation files were hashed before and
after and are unchanged (checked $=30$, changed $=0$). Only test files changed.

\paragraph{Four real blind spots.} The surviving mutants exposed four genuine
blind spots, all at the same structural location: consistency within a single
write batch. Two concern supersession---deleting the working-copy invalidation
leaves two tombstones for one old entry when two new values replace it in one
batch, and reversing tombstone order lets a superseded value return to
\emph{active} under last-write-wins in a chained batch. One concerns conflict
detection: its active-status filter is a defensive layer that end-to-end tests
can never reach, because retrieval has already filtered, so only a direct unit
test can pin it. One concerns the quality-label vocabulary: every label test
sentence also matched a second keyword, so deleting a keyword left the suite
green; the fix pins each label to a single keyword and asserts that the quoted
evidence is a verbatim substring. Each was reproduced, pinned by a minimal
regression assertion verified to pass on the original code and fail on the
mutant, and added to the repository (22 Python, 15 JavaScript). The two remaining
survivors are equivalent mutants: the negative
branch of the integral clamp is unreachable under the invariant
$\mathit{integral}\ge 0$ that holds for the contract's positive gain, so
$42/42=100\%$ of \emph{non-equivalent} mutants are killed. We report exactly
that sentence, not a blanket claim of full detection.

\paragraph{Assertion strength.} A fixed-seed sample of 238 assertions (165
Python, 73 JavaScript) was graded heuristically: 179 strong
($75.2\%$), 59 weak ($24.8\%$), and zero tautologies in the sample. A full
heuristic scan of 1{,}192 static assertions rated Python at 597 strong, 227
weak, and one tautology, and JavaScript at 278 strong and 89 weak, i.e.\
$26.6\%$ weak-or-tautological overall. We found and fixed one tautological
assertion (an import already performed by the module header), two
oracle-self-referential assertions whose expected values were read from the
implementation's own unsafe-control-action table (rewritten against a test-side
literal table; that file grew from 159 to 161 assertions), and one assertion
whose threshold inputs were taken from the audited configuration. Every rewrite
was validated by injecting the corresponding fault and observing the new
assertion fail. We therefore do not claim that 892 assertions are 892
independent precise literals. We claim that the sample contains no tautology,
that the full scan found one (now fixed), and that roughly a quarter of
assertions are graded weak.

\paragraph{Exception paths.} Under a strict criterion that counts only invalid,
empty, oversized, injected, timed-out, malformed, or boundary inputs, the
exception-path share is $199/825=24.1\%$ for Python and $82/367=22.3\%$ for
JavaScript (relaxed upper bounds: $43.3\%$ and $54.2\%$). Our own bar for
demanding new robustness tests was $20\%$; both suites clear it, but barely, and
excluding threshold-boundary assertions would take Python to about $11\%$. The
accurate reading is ``just over the line, not ample''.

\paragraph{Coverage.} Coverage was measured, not estimated: PEP~669
\texttt{sys.monitoring} for Python and V8 coverage for JavaScript. For
\texttt{engine/}, line coverage is $89.95\%$ and branch coverage $80.71\%$;
\texttt{control.py} reaches $141/141$ lines and $56/56$ branches. Whole-repository
Python coverage is $45.81\%$ line and $43.17\%$ branch, pulled down by
\texttt{tools/} at $32.92\%$, where only 3 of 26 files execute. We do not
use coverage to argue the absence of dead code: under the tested configuration
\texttt{control.py} has no unexecuted line, while two equivalent mutants
survive---which is why coverage alone cannot establish test strength.

\paragraph{Baseline bookkeeping.} The audited baseline was 13 Python files with
868 assertions and 381 JavaScript assertions; the audit added 22 and 2 Python
assertions and 15 JavaScript assertions, for 14 files with 892 and 396
respectively, all passing, with the implementation source unchanged.

\paragraph{Limits.} Four limits bound these claims. The repository has no version
control history, so we can characterise the current state but cannot prove that
no assertion was ever weakened. Mutation was single-fault only and did not cover
every module. No cross-language differential test exists: no mutant was killed
by both suites, so one-sided semantic drift is visible only to its own suite.
Finally, \texttt{tools/} coverage is a lower bound: subprocess-invoked tools are
recorded as unexecuted.
